\section{最佳逼近}

所谓逼近，就是用简单函数去近似复杂函数。数学分析中的 Weierstrass 第一逼近定理用 Bernstein 多项式具体构造了多项式序列，它可以在任何闭区间上一致逼近连续函数。实变函数中，可测函数可以用简单函数逼近……逼近是一个朴素的思想。

评价逼近好坏有不同的标准：逼近速度、精度、平滑程度等。为此，需要定量研究逼近误差。下面介绍目标函数、逼近工具和误差尺度等基本概念，并证明最佳逼近的存在性、唯一性与特征刻画。

\subsection{准备工作与基础}
先约定一些记号。
\begin{definition}
\begin{enumerate}
    \item $C[a,b]$ 表示闭区间 $[a,b]$ 上的连续函数全体；
    \item $C_{2\pi}$ 表示 $\R$ 上周期为 $2\pi$ 的连续函数全体；
    \item $P_n$ 为次数不超过 $n$ 的实多项式全体
    $\{a_nx^n+\cdots+a_1x+a_0: a_i\in\R\}$；
    \item $T_n$ 为次数不超过 $n$ 的三角多项式全体
    $\{a_n\cos nx+b_n\sin nx+\cdots +a_1\cos x+b_1\sin x +a_0: a_i,b_i\in\R\}$；
    \item 约定用 $p$ 表示多项式，$t$ 表示三角多项式；
    \item $\|f\| = \max_{x\in D} |f(x)|$，其中 $D$ 为函数的定义域。
\end{enumerate}
\end{definition}

经典的 Weierstrass 第一逼近定理表述为：
\begin{theorem}
设 $f\in C[a,b]$，则对任意 $\varepsilon>0$，存在多项式 $p$ 使得 $|f(x)-p(x)|<\varepsilon$ 对一切 $x\in[a,b]$ 成立。
\end{theorem}
该定理可通过构造 Bernstein 多项式
\[
B_n^f(x) = \sum_{k=0}^{n} f\left(\frac{k}{n}\right)\binom{n}{k}x^k(1-x)^{n-k}
\]
并借助平移伸缩推广到一般区间。

关于三角级数，对周期 $2\pi$ 的连续函数 $f$，其 Fourier 级数的系数为
\[
a_n = \frac{1}{\pi}\int_{-\pi}^{\pi}f(x)\cos nx\,dx\;(n\ge0),\qquad
b_n = \frac{1}{\pi}\int_{-\pi}^{\pi}f(x)\sin nx\,dx\;(n\ge1).
\]
记部分和
\[
S_n(x) = \frac{a_0}{2} + \sum_{k=1}^{n}(a_k\cos kx + b_k\sin kx),
\]
Cesàro 和
\[
F_n(x) = \frac{1}{n}\sum_{k=0}^{n-1} S_k(x),
\]
延迟平均
\[
V_n(x) = \frac{1}{n}\sum_{k=n}^{2n-1} S_k(x).
\]
连续函数的 Fourier 级数不一定处处收敛（卡尔森证明了几乎处处收敛），但其 Cesàro 和是一致收敛的，这也给出了 Weierstrass 第二逼近定理的构造性证明：
\begin{theorem}
设 $f\in C_{2\pi}$，则对任意 $\varepsilon>0$，存在三角多项式 $t$ 使得 $|f(x)-t(x)|<\varepsilon$ 对一切 $x\in\R$ 成立。
\end{theorem}
\begin{theorem}
设 $f\in C_{2\pi}$，则 $F_n(x)$ 一致收敛于 $f(x)$。
\end{theorem}

\subsection{误差尺度的衡量}
用次数不超过 $n$ 的多项式（或三角多项式）逼近 $f$，自然关心能达到的最小误差。定义
\[
e(f,P_n) = \inf\{\|f-p\| : p\in P_n\},
\]
以及 $e(f,T_n)$ 类似。易见对任意 $n$，存在子列使得该误差趋于零，但收敛速度可能很慢（例如 $f(x)=|x|$ 在 $0$ 附近的逼近）。因此需要在限定次数的条件下寻找“最佳逼近”。

\subsection{最佳逼近的存在性}
\begin{definition}
称 $p\in P_n$ 为 $f$ 在 $P_n$ 中的\textbf{最佳逼近}，若 $\|f-p\| = e(f,P_n)$。类似定义三角多项式的最佳逼近。
\end{definition}

\begin{theorem}
设 $f\in C[a,b]$，则存在 $p\in P_n$ 使 $\|f-p\| = e(f,P_n)$。
\end{theorem}
\begin{proof}
先给出一个引理。
\begin{lemma}
设 $\{p_m\}$ 是次数不超过 $n$ 的有界多项式列，则存在收敛子列 $p_{m_k}$ 满足 $\lim_{k\to\infty} p_{m_k} = p^*$。
\end{lemma}
引理可由反复使用 Bolzano–Weierstrass 定理证明。

下记 $e_n(f)=e(f,P_n)$。对每个 $m\ge1$，存在 $p_m\in P_n$ 使得
\[
\|f-p_m\|  \leqslant  e_n(f) + \frac1m  \leqslant  e_n(f)+1.
\]
因为 $e_n(f) \leqslant  \|f-0\| = \|f\|$，故
\[
\|p_m\|  \leqslant  \|f\| + \|p_m-f\|  \leqslant \|f\| + e_n(f)+1  \leqslant 2\|f\|+1.
\]
即 $\{p_m\}$ 有界，从而存在子列 $\{p_{m_k}\}$ 收敛到某个 $p^*\in P_n$。
下证 $p^*$ 即为最佳逼近。因为
\[
\|f-p_{m_k}\|  \leqslant e_n(f) + \frac1{m_k}  \leqslant e_n(f) + \frac1k,
\]
令 $k\to\infty$ 得 $\|f-p^*\| \leqslant e_n(f)$。又因 $e_n(f)$ 是下确界，必有 $\|f-p^*\| = e_n(f)$。
\end{proof}

类似可证
\begin{theorem}
设 $f\in C_{2\pi}$，则存在 $t\in T_n$ 使 $\|f-t\| = e(f,T_n)$。
\end{theorem}

\subsection{最佳逼近的特征刻画——切比雪夫交错定理}
\begin{definition}[交错偏离点]
设 $f\in C[a,b]$，$p\in P_n$，$r(x)=f(x)-p(x)$，
若存在点列
\[
a \leqslant  x_1 < x_2 < \dots < x_N  \leqslant  b,
\]
满足：
\begin{enumerate}
    \item 对每个 $j$，$|r(x_j)| = \|r\| $；
    \item 对每个 $j=1,\dots,N-1$，$r(x_j) = -r(x_{j+1})$，
\end{enumerate}
则称 $\{x_1,\dots,x_N\}$ 为 $r$ 的一组\textbf{交错点}，每个 $x_j$ 称为\textbf{交错偏离点}。
\end{definition}

\begin{theorem}[切比雪夫最佳逼近定理]
设 $f\in C[a,b]$，则 $p\in P_n$ 是 $f$ 的最佳逼近当且仅当误差函数 $r(x)=f(x)-p(x)$ 存在 $n+2$ 个交错偏离点。\\
设 $f\in C_{2\pi}$，则 $t\in T_n$ 是 $f$ 的最佳逼近当且仅当误差函数 $r(x)=f(x)-t(x)$ 存在 $2n+2$ 个交错偏离点。
\end{theorem}

\begin{proof}
仅对多项式情形给出证明，三角情形类似。

\textbf{充分性}：设 $x_1<\cdots<x_{n+2}$ 是 $p$ 的一组交错偏离点，即
\[
|f(x_j)-p(x_j)| = e_n(f),\quad j=1,\dots,n+2,
\]
且相邻点误差符号交替。不妨设 $f(x_j)-p(x_j) = (-1)^{j-1}e_n(f)$。
若 $p$ 不是最佳逼近，则存在 $q\in P_n$ 使 $\|f-q\| < \|f-p\|$。
考虑 $d(x)=p(x)-q(x)\in P_n$。在 $x_j$ 处，
\[
d(x_j) = (p(x_j)-f(x_j)) - (q(x_j)-f(x_j)).
\]
因 $\|f-q\|<\|f-p\|$，故 $d(x_j)$ 的符号与 $p(x_j)-f(x_j)$ 相同。于是 $d$ 在 $x_j$ 处依次变号，从而 $d$ 至少有 $n+1$ 个零点，但 $d$ 是 $n$ 次多项式，于是 $d$ 至少有$ n+1$ 个零点。
但 $d$ 为次数不超过 $n$ 的多项式，
除非 $d\equiv0$，
否则不可能有 $n+1$ 个不同零点。
故 $d\equiv0$，即 $p\equiv q$，
与假设矛盾。

\textbf{必要性}：
\begin{enumerate}
    \item \textbf{交错偏离点有限且小于 $n+2$} 

若交错偏离点个数 $N < n+2$，但 $p$ 仍是最佳逼近。记 $r(x)=f(x)-p(x)$，$M=\|r\|$。设
\[
a<x_1<\cdots<x_N<b
\]
为全体交错偏离点。

对每个 $i=1,\dots,N-1$，
由连续性可知存在
\[
\xi_i\in(x_i,x_{i+1})
\]
使得
\[
r(\xi_i)=0.
\]

再记
\[
\xi_0=a,\qquad \xi_N=b.
\]将 $[a,b]$ 用 $$\xi_1,\dots,\xi_{N-1}$$ 分割成若干子区间，则 $r$ 在每个子区间上满足
\[
-M+\alpha < r(x)  \leqslant M \quad\text{或}\quad -M  \leqslant r(x) < M-\alpha.
\]
因为我们可以假设在区间 $[a, \xi_1]$ 上有 $r(x)  \leqslant 0$。因为 $\|r\| = M$，从而 $|r(x)|  \leqslant  M$，从而此时 $r(x) \geqslant -M$ 是成立的。

由闭区间上连续函数的性质，存在 $0 < \alpha_1 < M $ 使得 $-M < r(x) < M -\alpha_1$ 在 $[a, \xi_1]$ 上成立。同理在每个子区间上都存在 $ 0 < \alpha_i < M $ 使得上述不等式成立。取 $\alpha = \min\{\alpha_i\}$，则在每个子区间上都有
\[-M   \leqslant  r(x) < M - \alpha.
\]或者
\[-M + \alpha < r(x)  \leqslant  M \]成立。

由于 $r$ 在每个子区间上不变号，
且 $\phi$ 在穿过每个 $\xi_i$ 后也恰变号一次，
故 $r$ 与 $\phi$ 在每个子区间上同号。

令 $\phi(x) = \prod_{i=1}^{N-1}(x-\xi_i)\in P_{N-1}\subseteq P_n$。构造 $q(x)=p(x)+c\phi(x)$。我们通过选取合适的 $c$ 去导出 $q(x)$ 是比 $p(x)$ 更好的逼近，从而导致矛盾。

不妨设在 $[a, \xi_1]$ 上 $r(x) \leqslant  0$ 并且 $\phi(x)$ 在 $[a, \xi_1]$ 上满足 $\phi(x)  \leqslant  0$（这只是符号上的问题）。
记 $K = \|\phi\|$.

此时 $-M  \leqslant  r(x) < M-\alpha $  注意到 $\phi(x)$ 与 $r(x)$ 在 $[a, \xi_1]$ 上符号相同。因为
$$
\begin{aligned}
f(x) - q(x) &= f(x) - p(x) - c\phi(x)  \\
              &= r(x) - c\phi(x) \\
              &= r(x) + c|\phi(x)| \\
\end{aligned}$$
这里由于 $r(x)  \leqslant  0$ 并且 $\phi(x)  \leqslant  0$，所以 
当 $c$ 满足 $ 0 < c| \phi(x) |  \leqslant  |c\cdot K| < \frac{\alpha}{2}$ 并且 $c>0$ 时
有
$$
\begin{aligned}
-M  < r(x) - c\phi(x) < M -\alpha + \frac{\alpha}{2} \\ 
\end{aligned}
$$
同理，在区间 $[\xi_1, \xi_2]$ 上我们能够利用
$$
\begin{aligned}
 -M +\alpha < r(x)  \leqslant  M
\end{aligned}
$$
和
$$
\begin{aligned}
f(x) - q(x) &= f(x) - p(x) - c\phi(x)  \\
              &= r(x) - c\phi(x) \\
              &= r(x) - c|\phi(x)| \\
\end{aligned}
$$
得到
$$
\begin{aligned}
-M + \alpha - \frac{\alpha}{2} < r(x) - c\phi(x) < M 
\end{aligned}$$
注意到上述过程与子区间的选择无关，因此在每个子区间上都满足上述不等式。
从而有 $$\|f-q\| < M = \|f-p\|$$，并且根据 $q$ 的构造可以得知它是不超过$n$ 次的多项式，
于是与 $\|f-q\| < \|f-p\|$矛盾。

\item \textbf{交错偏离点有限且大于等于 $n+2$} 

由于交错偏离点按符号交替出现，
故可从中依次选取
$n+2$
个点构成一组交错点。
由充分性部分，
$p$
即为最佳逼近。

\item \textbf{交错偏离点有无穷个}记若交错偏离点有无穷多个。记
$$
E_+=\{x\in[a,b]:r(x)=M\},\qquad
E_-=\{x\in[a,b]:r(x)=-M\}.
$$

由于 $r$ 连续，而 $\{M\},\{-M\}$ 为闭集，因此
$$
E_+=r^{-1}(\{M\}),\qquad
E_-=r^{-1}(\{-M\})
$$
均为闭集，从而为紧集。

由于交错偏离点无穷多个，因此在区间 $[a,b]$ 中，$E_+$ 与 $E_-$ 必无限次交替出现。

先取
$$
x_1=\min(E_+\cup E_-).
$$
不妨设
$$
r(x_1)=M.
$$

由于交错偏离点无穷多个，因此集合
$$
\{x>x_1:x\in E_-\}
$$
非空。又因为 $E_-$ 为闭集，所以该集合在 $[a,b]$ 中取到最小元。定义
$$
x_2=\min\{x>x_1:x\in E_-\}.
$$
于是
$$
x_1<x_2,\qquad r(x_2)=-M.
$$

同理，由于交错偏离点无穷多个，集合
$$
\{x>x_2:x\in E_+\}
$$
非空，并且取到最小元。定义
$$
x_3=\min\{x>x_2:x\in E_+\}.
$$
于是
$$
x_2<x_3,\qquad r(x_3)=M.
$$

继续递归构造：
若已定义
$$
x_k,
$$
则定义
$$
x_{k+1}
=
\min\{x>x_k:x\in E_{\sigma_{k+1}}\},
$$
其中
$$
\sigma_{k+1}=
\begin{cases}
+,& r(x_k)=-M,\\
-,& r(x_k)=M.
\end{cases}
$$

由于交错偏离点无穷多个，上述过程不会在有限步停止。

因此可得到
$$
x_1<x_2<\cdots<x_{n+2}
$$
满足
$$
r(x_i)=(-1)^iM.
$$

故已经存在不少于 $n+2$ 个交错偏离点。
\end{enumerate}
\end{proof}

\subsection{最佳逼近的唯一性}
\begin{theorem}
最佳逼近是唯一的：对 $f\in C[a,b]$，存在唯一的 $p\in P_n$ 使得 $\|f-p\| = e_n(f)$；对 $f\in C_{2\pi}$，存在唯一的 $t\in T_n$ 使得 $\|f-t\| = e(f,T_n)$。
\end{theorem}

\begin{proof}
若不然，设 $p,q\in P_n$ 均为最佳逼近，则 $\|f-p\| = \|f-q\| = e_n(f)$。由三角不等式，
\[
\left\|f - \frac{p+q}{2}\right\|  \leqslant \frac12\|f-p\| + \frac12\|f-q\| = e_n(f),
\]
故 $\frac{p+q}{2}$ 亦为最佳逼近。由切比雪夫交错定理，存在 $n+2$ 个交错偏离点 $x_1<\cdots<x_{n+2}$。记 $d_p = f-p$, $d_q = f-q$，则
\[
\left|f(x_1) - \frac{p(x_1)+q(x_1)}{2}\right| = \frac{|d_p(x_1)+d_q(x_1)|}{2} = e_n(f).
\]
因 $|d_p|,|d_q| \leqslant e_n(f)$，上式仅当 $d_p(x_1)=d_q(x_1)$ 时成立，从而 $p(x_1)=q(x_1)$。同理对所有 $x_i$ 均有 $p(x_i)=q(x_i)$。一个 $n$ 次多项式在 $n+1$ 个点取值相同必恒等，矛盾。
\end{proof}

利用唯一性可得一些有用的推论。
\begin{theorem}
设 $f\in C_{2\pi}$ 为偶函数，$t\in T_n$ 是其最佳逼近，则 $t$ 也是偶函数。
\end{theorem}
\begin{proof}
由 $\|f(x)-t(x)\| \leqslant e_n(f)$ 得 $\|f(-x)-t(-x)\| \leqslant e_n(f)$，故 $\|f(x)-t(-x)\| \leqslant e_n(f)$，即 $s(x)=t(-x)$ 也是最佳逼近。由唯一性 $t(x)=s(x)=t(-x)$。
\end{proof}